\chapter{Oppervlakken}\label{ch:b2:surfaces}

Na de krommen de oppervlakken: objecten met twee parameters in
$\R^3$. De differentiaalrekening van \cref{ch:b2:diffcalc} levert
alles wat wij nodig hebben --- partiële afgeleiden geven
raakvectoren, het kruisproduct geeft de normaal, \omterm{def:b2:linalg:det}{determinanten}
geven \omterm{def:b2:surfaces:area}{oppervlakten}. Wij definiëren reguliere \omterm{def:b2:surfaces:param}{geparametriseerde
oppervlakken}, hun raakvlakken en de \emph{\omterm{def:b2:surfaces:fff}{eerste
fundamentaalvorm}}, die alle lengte- en oppervlaktemetingen op het
oppervlak codeert. Oppervlakken duiken ook op als
niveauverzamelingen $f(x, y, z) = c$; de gradiënt richt dan de
normaal.

\section{Geparametriseerde oppervlakken}

\begin{definition}[Regulier geparametriseerd oppervlak]\label{def:b2:surfaces:param}
Zij $U \subseteq \R^2$ \omterm{def:b2:metric:topology}{open}. Een \emph{geparametriseerd
oppervlak}\index{geparametriseerd oppervlak} van klasse
$\mathcal{C}^k$ ($k \geq 1$) is een afbeelding $\sigma \colon U \to
\R^3$, $(u, v) \mapsto \sigma(u, v)$, van klasse $\mathcal{C}^k$.
Een punt heet \emph{regulier} wanneer de vectoren van de partiële
afgeleiden
\[
\sigma_u(u, v) = \frac{\partial\sigma}{\partial u}(u,v),
\qquad
\sigma_v(u, v) = \frac{\partial\sigma}{\partial v}(u,v)
\]
lineair onafhankelijk zijn, dat wil zeggen $\sigma_u \wedge
\sigma_v \neq 0$; het oppervlak heet \emph{regulier} wanneer elk
punt dat is.
\end{definition}

\begin{example}[De drie standaardbeschrijvingen]\label{ex:b2:surfaces:standard}
\leavevmode
\begin{enumerate}
\item \emph{Grafiek:} $\sigma(u, v) = (u,\ v,\ f(u, v))$ voor $f
\in \mathcal{C}^1(U)$. Altijd regulier: $\sigma_u = (1, 0, f_u)$ en
$\sigma_v = (0, 1, f_v)$ zijn onafhankelijk.
\item \emph{Bol (bolcoördinaten):} voor de bol met straal $R$,
\[
\sigma(\theta, \varphi)
= (R\cos\theta\cos\varphi,\ R\sin\theta\cos\varphi,\
R\sin\varphi),
\qquad
(\theta, \varphi) \in \R \times
\bigl(-\tfrac\pi2, \tfrac\pi2\bigr),
\]
met $\theta$ de lengtegraad en $\varphi$ de breedtegraad. Men gaat
na dat $\norm{\sigma_\theta \wedge \sigma_\varphi} =
R^2\cos\varphi > 0$: regulier buiten de polen (die deze kaart
weglaat).
\item \emph{Niveauverzameling:} $S = \{(x,y,z) : f(x, y, z) = c\}$
met $f$ van klasse $\mathcal{C}^1$ en $\nabla f \neq 0$ op $S$.
Nabij elk punt kan één coördinaat als functie van de twee andere
worden uitgedrukt met de stelling van de impliciete functie
(\cref{ch:b2:diffcalc}), dus is $S$ lokaal een grafiek.
\end{enumerate}
\end{example}

\begin{example}[Van niveauverzameling naar grafiek]
\label{ex:b2:surfaces:levelgraph}
De stelling van de impliciete functie in punt 3 verdient één
expliciete doorloop. Neem de bol $x^2 + y^2 + z^2 = R^2$ nabij haar
noordpool $(0, 0, R)$: daar is $\frac{\partial f}{\partial z} = 2z
= 2R \neq 0$, en naar $z$ oplossen geeft de grafiekkaart
\[
z = \sqrt{R^2 - x^2 - y^2},
\qquad x^2 + y^2 < R^2 ,
\]
overal regulier op haar (\omterm{def:b2:metric:topology}{open}) domein --- ook in de pool die de
bolkaart miste. Nabij een punt op de evenaar zoals $(R, 0, 0)$ lost
dezelfde stelling naar $x$ op ($\frac{\partial f}{\partial x} = 2R
\neq 0$). De vuistregel: een niveauoppervlak is een grafiek boven
het coördinaatvlak dat \emph{loodrecht staat op de grootste
component van de gradiënt}, en door de bol met zes zulke
grafiekkaarten te overdekken, gaat men haar gladheid overal na
zonder enige goniometrie.
\end{example}

\section{Raakvlak en normaal}

\begin{definition}[Raakvlak]\label{def:b2:surfaces:tangent}
Zij $\sigma$ regulier in $(u_0, v_0)$ en $M_0 = \sigma(u_0, v_0)$.
Het \emph{raakvlak}\index{raakvlak} $T_{M_0}S$ is het vlak door
$M_0$ met richting $\operatorname{Vect}(\sigma_u, \sigma_v)$ (de
partiële afgeleiden in $(u_0, v_0)$). De
\emph{eenheidsnormaal}\index{eenheidsnormaal} is
\[
n(u_0, v_0)
= \frac{\sigma_u \wedge \sigma_v}
       {\norm{\sigma_u \wedge \sigma_v}} .
\]
\end{definition}

\begin{proposition}[Raakvectoren zijn
snelheidsvectoren]\label{prop:b2:surfaces:velocity}
De richting van $T_{M_0}S$ is precies de verzameling vectoren
$\gamma'(0)$, waarbij $\gamma = \sigma \circ c$ de
$\mathcal{C}^1$-krommen doorloopt die op het oppervlak door $M_0$
zijn getekend (dat wil zeggen: $c \colon (-\varepsilon,
\varepsilon) \to U$ is $\mathcal{C}^1$ met $c(0) = (u_0, v_0)$).
\end{proposition}

\begin{proof}
Is $c(t) = (u(t), v(t))$, dan geeft de kettingregel
(\cref{ch:b2:diffcalc})
\[
\gamma'(0) = u'(0)\,\sigma_u + v'(0)\,\sigma_v
\in \operatorname{Vect}(\sigma_u, \sigma_v).
\]
Omgekeerd wordt de vector $a\sigma_u + b\sigma_v$ bereikt door de
kromme $c(t) = (u_0 + at,\ v_0 + bt)$, die voor kleine $\abs t$ in
de \omterm{def:b2:metric:topology}{open} verzameling $U$ blijft.
\end{proof}

\begin{example}[Het raakvlak van de helicoïde]
\label{ex:b2:surfaces:helicoidtangent}
Voor de helicoïde $\sigma(u, v) = (v\cos u,\ v\sin u,\ au)$ is in
het punt $\sigma(0, 1) = (1, 0, 0)$:
\[
\sigma_u = (0,\ 1,\ a), \qquad \sigma_v = (1,\ 0,\ 0),
\qquad
\sigma_u \wedge \sigma_v = (0,\ a,\ -1),
\]
dus is het \omterm{def:b2:surfaces:tangent}{raakvlak} $a\,y = z$. Het bevat de hele horizontale
beschrijvende rechte $\{(t, 0, 0)\}$ (richting $\sigma_v$): net als
bij de kegel van \cref{exo:b2:surfaces:1} ligt bij een oppervlak
dat door rechten wordt beschreven, elke beschrijvende rechte in het
\omterm{def:b2:surfaces:tangent}{raakvlak} eraan. De andere raakrichting $\sigma_u$ is de snelheid
van de schroeflijn $u \mapsto \sigma(u, 1)$: één kaart, twee
getekende krommen, en het hele \omterm{def:b2:surfaces:tangent}{raakvlak} is opgespannen ---
\cref{prop:b2:surfaces:velocity} in actie.
\end{example}

\begin{proposition}[Normaal van een niveauoppervlak]\label{prop:b2:surfaces:gradient}
Zij $S = \{f = c\}$ met $f$ van klasse $\mathcal{C}^1$ en $\nabla
f(M_0) \neq 0$. Dan is het \omterm{def:b2:surfaces:tangent}{raakvlak} van $S$ in $M_0$ het vlak door
$M_0$ loodrecht op $\nabla f(M_0)$:
\[
T_{M_0}S :\quad
\langle \nabla f(M_0),\ M - M_0\rangle = 0 .
\]
\end{proposition}

\begin{proof}
Voor elke kromme $\gamma$ getekend op $S$ door $M_0$ geldt
identiek $f(\gamma(t)) = c$, dus geeft de kettingregel dat $\langle
\nabla f(M_0), \gamma'(0)\rangle = 0$: alle snelheidsvectoren staan
loodrecht op de gradiënt, dus de raakrichting is bevat in het vlak
$\nabla f(M_0)^\perp$. Beide zijn tweedimensionale deelruimten ---
de raakrichting omdat $S$ lokaal een reguliere grafiek is
(\cref{ex:b2:surfaces:standard}), het orthogonale complement omdat
$\nabla f(M_0) \neq 0$ --- en dus zijn zij gelijk.
\end{proof}

\begin{example}\label{ex:b2:surfaces:spheretangent}
Voor de bol $x^2 + y^2 + z^2 = R^2$ is $\nabla f = 2(x, y, z)$, dus
staat het \omterm{def:b2:surfaces:tangent}{raakvlak} in $M_0$ loodrecht op de straal $\vect{OM_0}$
--- het klassieke feit dat straal en \omterm{def:b2:surfaces:tangent}{raakvlak} loodrecht op elkaar
staan, met vergelijking $\langle M_0, M\rangle = R^2$.
\end{example}

\begin{example}[Raakvlak van een grafiek]\label{ex:b2:surfaces:graphtangent}
Voor $z = f(x, y)$ in $(x_0, y_0)$ geeft
\cref{prop:b2:surfaces:gradient} toegepast op $F(x,y,z) = f(x,y) -
z$:
\[
z = f(x_0, y_0)
+ f_x(x_0, y_0)(x - x_0)
+ f_y(x_0, y_0)(y - y_0) ,
\]
het \omterm{def:b2:affine:subspace}{affiene} deel van de eerste-ordeontwikkeling van Taylor --- het
\omterm{def:b2:surfaces:tangent}{raakvlak} is de grafiek van de \omterm{def:b2:diffcalc:differential}{differentiaal}, zoals het moet.
\end{example}

\begin{example}[Het dichtstbijzijnde punt van een
oppervlak]\label{ex:b2:surfaces:closest}
Welk punt van de paraboloïde $z = x^2 + y^2$ ligt het dichtst bij
$P = (0, 0, 1)$? Minimaliseer de kwadratische afstand langs het
oppervlak: met $\rho^2 = x^2 + y^2$ is
\[
g(\rho^2) = \rho^2 + (\rho^2 - 1)^2,
\qquad
g'(\rho^2) = 1 + 2(\rho^2 - 1) = 0
\iff \rho^2 = \tfrac12 ,
\]
wat de cirkel van punten op hoogte $z = \frac12$ en afstand
$\sqrt{\tfrac12 + \tfrac14} = \frac{\sqrt3}2$ geeft. De meetkundige
handtekening van de minimaliteit: in zulk een punt $M$ moet de
vector $\vect{MP}$ \emph{normaal} op het oppervlak staan --- anders
zou langs een getekende kromme glijden met een snelheid die een
component naar $P$ heeft, de afstand verkleinen. Controle:
$\nabla(z - x^2 - y^2) = (-2x, -2y, 1)$ in $M = (x, y, \tfrac12)$,
terwijl $\vect{MP} = (-x, -y, \tfrac12) = \tfrac12(-2x, -2y, 1)$:
evenwijdig, zoals voorspeld. De eerste-ordevoorwaarde ``voetpunt
van de loodlijn'' is dezelfde die in \cref{exo:b2:surfaces:12} de
extrema op niveauverzamelingen zal aandrijven.
\end{example}

\begin{example}[Raakvlakken van kwadrieken: de
polarisatieregel]\label{ex:b2:surfaces:polarization}
Zij $S : xy + yz + zx = 1$ en $M_0 = (1, 1, 0) \in S$. Hier is
$\nabla f = (y + z,\ x + z,\ x + y)$, dus $\nabla f(M_0) = (1, 1,
2)$ en is het \omterm{def:b2:surfaces:tangent}{raakvlak}
\[
(x - 1) + (y - 1) + 2z = 0,
\qquad\text{dat wil zeggen}\qquad x + y + 2z = 2 .
\]
Hetzelfde antwoord komt uit de \emph{polarisatieregel} die
\cref{ex:b2:surfaces:spheretangent} en \cref{exo:b2:surfaces:2}
veralgemeent: vervang in de vergelijking van de \omterm{pb:b2:surfaces:1}{kwadriek} $x^2$ door
$x_0x$, en elk product $xy$ door $\frac{x_0y + y_0x}2$ (en
cyclisch):
\[
\frac{x_0y + y_0x}2 + \frac{y_0z + z_0y}2 + \frac{z_0x +
x_0z}2 = 1
\;\xrightarrow{\ M_0 = (1,1,0)\ }\;
\frac{x + y}2 + \frac z2 + \frac z2 = 1 ,
\]
wat opnieuw $x + y + 2z = 2$ is. De regel werkt omdat $\nabla$ van
een \omterm{def:b2:quadratic:def}{kwadratische vorm} de bijbehorende bilineaire vorm is,
geëvalueerd tegen het basispunt --- raken aan een \omterm{pb:b2:surfaces:1}{kwadriek} is
polarisatie, nog een gezicht van \cref{ch:b2:quadratic}.
\end{example}

\section{De eerste fundamentaalvorm}

\begin{definition}[Eerste fundamentaalvorm]\label{def:b2:surfaces:fff}
Zij $\sigma \colon U \to \R^3$ een regulier
$\mathcal{C}^1$-oppervlak. Haar \emph{eerste
fundamentaalvorm}\index{eerste fundamentaalvorm} in $(u,v)$ is de
positief definiete \omterm{def:b2:quadratic:def}{kwadratische vorm} op $\R^2$
\[
I(h, k) = \norm{h\,\sigma_u + k\,\sigma_v}^2
= E\,h^2 + 2F\,hk + G\,k^2,
\]
waarbij
\[
E = \norm{\sigma_u}^2,
\qquad
F = \langle \sigma_u, \sigma_v\rangle,
\qquad
G = \norm{\sigma_v}^2 .
\]
\end{definition}

\begin{remark}
$I$ is de beperking van het omgevende euclidische inproduct tot het
\omterm{def:b2:surfaces:tangent}{raakvlak}, gelezen in de basis $(\sigma_u, \sigma_v)$: zij is
positief definiet precies omdat $\sigma_u, \sigma_v$ onafhankelijk
zijn (\cref{ch:b2:quadratic}). Elke metrische grootheid op het
oppervlak --- \omterm{def:b2:curves:length}{lengten} van getekende krommen, hoeken ertussen,
\omterm{def:b2:surfaces:area}{oppervlakten} --- wordt uit $E, F, G$ alleen berekend. Twee
oppervlakken met dezelfde $E, F, G$ in geschikte parameters zijn
\emph{isometrisch}, ook al liggen zij verschillend in de ruimte:
dat is het startpunt van de intrinsieke meetkunde.
\end{remark}

\begin{example}[Hoeken tussen
coördinaatkrommen]\label{ex:b2:surfaces:coordangle}
De \omterm{def:b2:surfaces:fff}{eerste fundamentaalvorm} meet ook hoeken: de coördinaatkrommen $u
\mapsto \sigma(u, v_0)$ en $v \mapsto \sigma(u_0, v)$ ontmoeten
elkaar onder de hoek $\theta$ met
\[
\cos\theta = \frac{\langle\sigma_u,
\sigma_v\rangle}{\norm{\sigma_u}\,\norm{\sigma_v}}
= \frac{F}{\sqrt{EG}} :
\]
de ene coëfficiënt $F$ beslist over de orthogonaliteit van het
parameternet. Voor de bolkaart en de helicoïde is $F = 0$:
meridianen snijden parallellen, en schroeflijnen snijden de
horizontale beschrijvende rechten, loodrecht --- en daarom klapten
hun oppervlakte-integranden ineen tot $\sqrt{EG}$. Voor een
grafiekkaart is $F = f_xf_y$, dat alleen verdwijnt waar een
partiële afgeleide dat doet: het coördinatennet van een gekantelde
grafiek is \emph{niet} orthogonaal, ook al is het $(x, y)$-net
eronder dat wel. Zien berekeningen op een oppervlak er zwaar uit,
dan is de eerste zet: zoek een kaart met $F = 0$.
\end{example}

\begin{example}[De zadelkaart]\label{ex:b2:surfaces:saddlechart}
Voor het zadel $z = xy$ met kaart $\sigma(u, v) = (u, v, uv)$:
\[
\sigma_u = (1, 0, v), \qquad \sigma_v = (0, 1, u),
\qquad
E = 1 + v^2, \quad F = uv, \quad G = 1 + u^2 ,
\]
en $EG - F^2 = 1 + u^2 + v^2 > 0$: overal regulier. De twee
coördinaatkrommen door een punt zijn \emph{rechte lijnen} van
$\R^3$ (houd $u$ of $v$ vast: de beschrijvende rechten van het
dubbel beschreven zadel), en toch is $F \neq 0$ buiten de assen: de
beschrijvende rechten door een generiek punt staan niet loodrecht
op elkaar. Beide beschrijvende rechten liggen in het \omterm{def:b2:surfaces:tangent}{raakvlak}, dat
zij opspannen --- dus snijdt het \omterm{def:b2:surfaces:tangent}{raakvlak} het oppervlak langs twee
hele rechten, het uiterste tegendeel van de bol, waarvan de
raakvlakken in slechts één punt raken. Het teken van het
``contact van tweede orde'' tussen een oppervlak en zijn
raakvlakken is een verhaal over \omterm{thm:b2:curves:frenet2d}{kromming}, dat in het volume van
bachelorjaar 3 wordt opgenomen.
\end{example}

\begin{proposition}[Lengte van een op een oppervlak getekende
kromme]\label{prop:b2:surfaces:curvelength}
Is $\gamma(t) = \sigma(u(t), v(t))$, $t \in [a, b]$, van klasse
$\mathcal{C}^1$, dan is
\[
L(\gamma) = \int_a^b
\sqrt{E\,u'^2 + 2F\,u'v' + G\,v'^2}\;\dd t ,
\]
met $E, F, G$ geëvalueerd in $(u(t), v(t))$.
\end{proposition}

\begin{proof}
$\gamma' = u'\sigma_u + v'\sigma_v$ volgens de kettingregel, dus
$\norm{\gamma'}^2 = I(u', v') = Eu'^2 + 2Fu'v' + Gv'^2$; integreer
$\norm{\gamma'}$ (\cref{def:b2:curves:length}).
\end{proof}

\begin{example}[Waarom lijnvliegtuigen over de pool vliegen]\label{ex:b2:surfaces:polarroute}
Twee luchthavens liggen op breedtegraad $\varphi_0$ en op
tegengestelde lengtegraden: $A = \sigma(0, \varphi_0)$ en $B =
\sigma(\pi, \varphi_0)$ op de bol met straal $R$. Langs de parallel
($\varphi \equiv \varphi_0$) is de \omterm{def:b2:curves:length}{lengte} $\int_0^\pi
R\cos\varphi_0\,\dd\theta = \pi R\cos\varphi_0$. Langs de route
over de pool (langs de meridiaan $\theta = 0$ omhoog en langs de
meridiaan $\theta = \pi$ omlaag) is zij $2R(\frac\pi2 -
\varphi_0)$. Op breedtegraad $\varphi_0 = \frac\pi3$ (zestig
graden): route over de parallel $\pi R/2 \approx 1.571\,R$, route
over de pool $\pi R/3 \approx 1.047\,R$ --- een derde korter.
Sterker nog, $\pi\cos\varphi_0 \geq \pi - 2\varphi_0$ op
$\intcc0{\pi/2}$ (de functie $\pi\cos\varphi - \pi + 2\varphi$
verdwijnt in beide uiteinden en haar afgeleide $2 - \pi\sin\varphi$
verandert eenmaal van teken, dus is zij eerst stijgend en dan
dalend, en bijgevolg niet-negatief): de polaire route verliest
nooit. De \omterm{def:b2:surfaces:fff}{eerste fundamentaalvorm} maakte van een navigatievraag
twee integralen van één regel; \cref{exo:b2:surfaces:6} drijft het
idee door tot een echt minimaliteitsbewijs voor meridianen.
\end{example}

\begin{lemma}[Identiteit van Lagrange]\label{lem:b2:surfaces:lagrange}
Voor alle $a, b \in \R^3$ geldt $\norm{a \wedge b}^2 = \norm a^2
\norm b^2 - \langle a, b\rangle^2$. In het bijzonder is
\[
\norm{\sigma_u \wedge \sigma_v} = \sqrt{EG - F^2}.
\]
\end{lemma}

\begin{proof}
Beide leden veranderen niet als wij $b$ vervangen door haar
component $b_\perp = b - \frac{\langle a, b\rangle}{\norm a^2}a$
loodrecht op $a$ (voor $a \neq 0$; het geval $a = 0$ is triviaal):
het linkerlid omdat $a \wedge a = 0$, het rechterlid door
$\norm{b_\perp}^2 = \norm b^2 - \frac{\langle a,b\rangle^2}{\norm
a^2}$ uit te werken en $\langle a, b_\perp\rangle = 0$. Het
volstaat dus de identiteit voor orthogonale $a, b$ te bewijzen,
waar zij luidt $\norm{a \wedge b} = \norm a \norm b$: waar, want
voor orthogonale vectoren heeft het kruisproduct \omterm{def:b2:nvs:norm}{norm} $\norm
a\norm b\,\abs{\sin\frac\pi2}$. De getoonde formule is het geval $a
= \sigma_u$, $b = \sigma_v$.
\end{proof}

\begin{remark}
De identiteit van Lagrange zegt dat $EG - F^2 =
\det\begin{psmallmatrix} E & F\\ F & G\end{psmallmatrix}$ de
\emph{\omterm{def:b2:linalg:det}{Gram-determinant}} van $(\sigma_u, \sigma_v)$ is: het kwadraat
van de \omterm{def:b2:surfaces:area}{oppervlakte} van het parallellogram dat zij opspannen.
Regulariteit, positieve definietheid van de \omterm{def:b2:surfaces:fff}{eerste
fundamentaalvorm} en positiviteit van de \omterm{def:b2:linalg:det}{Gram-determinant} zijn drie
formuleringen van één voorwaarde --- en daarom verdwijnt de
oppervlakte-integrand hieronder nooit op een reguliere kaart.
\end{remark}

\begin{definition}[Oppervlakte]\label{def:b2:surfaces:area}
Zij $\sigma \colon U \to \R^3$ een regulier injectief
$\mathcal{C}^1$-oppervlak en $K \subseteq U$ een \omterm{def:b2:metric:compact}{compact} gebied
waarop dubbele integralen zin hebben (\cref{ch:b2:multint}). De
\emph{oppervlakte}\index{oppervlakte!van een oppervlak} van het
stuk $\sigma(K)$ is
\[
\mathcal{A}
= \iint_K \norm{\sigma_u \wedge \sigma_v}\,\dd u\,\dd v
= \iint_K \sqrt{EG - F^2}\;\dd u\,\dd v .
\]
\end{definition}

\begin{remark}[Waarom deze formule]
De rechthoek $[u, u + \dd u] \times [v, v + \dd v]$ wordt tot op
eerste orde afgebeeld op het parallellogram opgespannen door
$\sigma_u\,\dd u$ en $\sigma_v\,\dd v$, waarvan de \omterm{def:b2:surfaces:area}{oppervlakte}
$\norm{\sigma_u \wedge \sigma_v}\,\dd u\,\dd v$ is: de definitie
integreert de lokale factor van oppervlaktevervorming, precies
zoals de \omterm{def:b2:curves:length}{booglengte} de lokale snelheid integreert. De
verenigbaarheid met \omterm{def:b2:curves:reparam}{parameterveranderingen} is
\cref{exo:b2:surfaces:8}; de verenigbaarheid met de formule voor de
verandering van veranderlijken bij dubbele integralen wordt in
\cref{ch:b2:multint} besproken.
\end{remark}

\begin{example}[Twee verschillende grafieken, één
oppervlakte]\label{ex:b2:surfaces:samearea}
Vergelijk boven de eenheidsschijf de kom $z = \frac12(x^2 + y^2)$
met het zadel $z = xy$. Hun oppervlakte-integranden
(\cref{exo:b2:surfaces:5}) zijn
\[
\sqrt{1 + x^2 + y^2}
\qquad\text{en}\qquad
\sqrt{1 + y^2 + x^2} :
\]
identiek. De twee oppervlakken --- het ene in alle richtingen
dezelfde kant op gekromd, het andere zadelvormig --- hebben boven
elk gebied precies gelijke \omterm{def:b2:surfaces:area}{oppervlakten}, $\frac{2\pi}3(2\sqrt2 -
1)$ boven de eenheidsschijf. Het oppervlakte-element ziet alleen de
\emph{\omterm{def:b2:curves:length}{lengte}} van de gradiënt, niet de schikking van de buiging;
de kom van het zadel onderscheiden vereist gegevens van tweede orde
(de tekenstructuur uit \cref{fig:b2:surfaces:saddle}), die geen
enkele oppervlaktemeting opmerkt. De \omterm{def:b2:surfaces:fff}{eerste fundamentaalvorm}:
metrisch, blind voor de vorm; de tweede vorm, die de vorm wél ziet,
hoort bij bachelorjaar 3.
\end{example}

\begin{example}[Oppervlakte van de bol]\label{ex:b2:surfaces:spherearea}
Voor de bolkaart uit \cref{ex:b2:surfaces:standard}:
\[
\sigma_\theta = R(-\sin\theta\cos\varphi,\
\cos\theta\cos\varphi,\ 0),
\qquad
\sigma_\varphi = R(-\cos\theta\sin\varphi,\
-\sin\theta\sin\varphi,\ \cos\varphi),
\]
dus $E = R^2\cos^2\varphi$, $F = 0$, $G = R^2$ en $\sqrt{EG - F^2}
= R^2\cos\varphi$. Bijgevolg is
\[
\mathcal{A}
= \int_{-\pi/2}^{\pi/2}\!\!\int_0^{2\pi}
R^2\cos\varphi\;\dd\theta\,\dd\varphi
= 2\pi R^2\,\bigl[\sin\varphi\bigr]_{-\pi/2}^{\pi/2}
= \boxed{4\pi R^2} .
\]
\end{example}

\begin{example}[De kegel, getoetst aan de
schoolformule]\label{ex:b2:surfaces:conearea}
Voor de kegel $z = \sqrt{x^2 + y^2}$ boven de ring $a \leq \rho
\leq b$ geeft de grafiekformule van \cref{exo:b2:surfaces:5} dat $1
+ f_x^2 + f_y^2 = 2$ (bereken $f_x = x/\rho$, $f_y = y/\rho$), dus
\[
\mathcal A = \sqrt2\,\pi\,(b^2 - a^2) .
\]
Controle met de formule $\pi\rho\ell$ voor de schuine hoogte uit
\cref{ex:b2:surfaces:revolution}: de volledige kegels met
grondstralen $b$ en $a$ hebben manteloppervlakten $\pi b\cdot
b\sqrt2$ en $\pi a\cdot a\sqrt2$, waarvan het verschil precies
$\sqrt2\pi(b^2 - a^2)$ is. Twee kaarten, twee formules, één
\omterm{def:b2:surfaces:area}{oppervlakte} --- de invariantie bewezen in
\cref{exo:b2:surfaces:8}, in het wild waargenomen.
\end{example}

\begin{remark}[Verstandscontroles voor oppervlakten]
Drie onmiddellijke controles vangen de meeste fouten in een
oppervlakteberekening. \emph{Schalen}: een oppervlak met $\lambda$
oprekken vermenigvuldigt $E, F, G$ met $\lambda^2$ en de
\omterm{def:b2:surfaces:area}{oppervlakte} met $\lambda^2$ --- een antwoord waarvan de afhankelijkheid
van $R$ niet kwadratisch is (zoals $4\pi R^2$), is fout.
\emph{Positiviteit van het element}: $\sqrt{EG - F^2}$ moet strikt
positief zijn in het inwendige van de kaart; een verdwijnende
waarde signaleert een ontaarding van de kaart, die moet worden
weggesneden zoals bij de polen van de bol. \emph{Symmetrie}: een
berekening over een \omterm{def:b2:quadratic:adjoint}{symmetrisch} stuk moet verenigbaar zijn met het
optellen van haar congruente delen --- de halve bol moet wel
$2\pi R^2$ geven.
\end{remark}

\begin{example}[Omwentelingsoppervlak]\label{ex:b2:surfaces:revolution}
Draai de kromme $z \mapsto (r(z), 0, z)$, met $r > 0$ van klasse
$\mathcal{C}^1$, om de $z$-as:
\[
\sigma(\theta, z) = (r(z)\cos\theta,\ r(z)\sin\theta,\ z).
\]
Dan is $E = r(z)^2$, $F = 0$, $G = 1 + r'(z)^2$, dus
\[
\mathcal{A} = \int_a^b\!\!\int_0^{2\pi}
r(z)\sqrt{1 + r'(z)^2}\;\dd\theta\,\dd z
= 2\pi\int_a^b r(z)\sqrt{1 + r'(z)^2}\,\dd z ,
\]
de klassieke formule (omtrek $2\pi r$ maal het element van schuine
\omterm{def:b2:curves:length}{lengte}). Voor de kegel $r(z) = kz$, $z \in [0, h]$: $\mathcal A =
2\pi k\sqrt{1 + k^2}\,\frac{h^2}{2} = \pi \rho \ell$ met $\rho =
kh$ de grondstraal en $\ell = h\sqrt{1 + k^2}$ de schuine hoogte
--- de schoolformule, nu afgeleid in plaats van aangenomen.
\end{example}

\begin{example}[De catenoïde]\label{ex:b2:surfaces:catenoid}
Draai de kettinglijn $r(z) = \cosh z$, $z \in \intcc{-1}{1}$, om
haar as: de resulterende \emph{catenoïde} heeft, volgens de formule
voor omwentelingsoppervlakken en $1 + \sinh^2 z = \cosh^2 z$,
\[
\mathcal A = 2\pi\int_{-1}^1\cosh z\,\sqrt{1 + \sinh^2z}\,
\dd z = 2\pi\int_{-1}^1\cosh^2z\,\dd z
= \pi\bigl[z + \sinh z\cosh z\bigr]_{-1}^{1},
\]
dat wil zeggen
\[
\mathcal A = 2\pi + \pi\sinh 2 \approx 17.68 .
\]
De factor $\sqrt{1 + r'^2}$ voor de schuinte versmolt met de straal
tot een volmaakt kwadraat --- dezelfde identiteit die de \omterm{def:b2:curves:length}{booglengte}
van de kettinglijn in het hoofdstuk over krommen elementair maakte.
Dat is geen algebraïsch toeval: van alle omwentelingsoppervlakken
die de twee randcirkels opspannen, minimaliseert de catenoïde de
\omterm{def:b2:surfaces:area}{oppervlakte} (zij is de vorm van een zeepvlies tussen twee ringen),
en precies deze variationele eigenschap zondert $\cosh$ af; het
volume van bachelorjaar 3 bewijst het met de variatierekening.
\end{example}

\begin{figure}[ht]
\centering
\begin{tikzpicture}[scale=1.05]
  % A saddle surface z = x^2 - y^2 drawn as a wire mesh in oblique projection
  % projection: (x,y,z) -> (x + 0.45*y, z + 0.35*y)
  \begin{scope}
    % u-lines (x varies, y fixed)
    \foreach \yy in {-1,-0.5,0,0.5,1} {
      \draw[ocDarkBlue!70, thick, domain=-1:1, smooth, samples=25]
        plot ({\x + 0.45*\yy}, {\x*\x - \yy*\yy + 0.35*\yy});
    }
    % v-lines (y varies, x fixed)
    \foreach \xx in {-1,-0.5,0,0.5,1} {
      \draw[ocMint!80!black, thick, domain=-1:1, smooth, samples=25]
        plot ({\xx + 0.45*\x}, {\xx*\xx - \x*\x + 0.35*\x});
    }
    % tangent plane at origin: z = 0 -> parallelogram
    \draw[ocRed, dashed]
      (-0.8 - 0.36, -0.28) -- (0.8 - 0.36, -0.28)
      -- (0.8 + 0.36, 0.28) -- (-0.8 + 0.36, 0.28) -- cycle;
    \fill (0,0) circle (0.045) node[below right, font=\small] {$M_0$};
    % normal vector at origin
    \draw[->, thick] (0,0) -- (0,0.9)
      node[right, font=\small] {$n$};
  \end{scope}
\end{tikzpicture}
\caption{Het zadel $z = x^2 - y^2$ nabij de oorsprong, met zijn
coördinaatkrommen ($u$-krommen in het blauw, $v$-krommen in het
groen), het \omterm{def:b2:surfaces:tangent}{raakvlak} in $M_0 = (0,0,0)$ (gestreept) en de
\omterm{def:b2:surfaces:tangent}{eenheidsnormaal} $n$. Het oppervlak kruist zijn \omterm{def:b2:surfaces:tangent}{raakvlak} --- het
tweedimensionale analogon van een buigpunt.}
\label{fig:b2:surfaces:saddle}
\end{figure}

\begin{remark}[Klassieke valkuilen]
(i) \emph{Singulariteiten van de kaart zijn geen singulariteiten
van het oppervlak}: de bolkaart ontaardt in de polen ($\cos\varphi
= 0$), maar de bol is daar volmaakt glad --- een andere kaart
(verwissel de rollen van de assen) is regulier in de polen. Probeer
een tweede parametrisering voordat je een punt singulier noemt.
(ii) \emph{De regulariteit van $\sigma$ betreft de parametrisering,
niet het beeld}: $\sigma(u, v) = (u^3, v, 0)$ is niet regulier
langs $u = 0$, hoewel haar beeld een vlak is. (iii) De
oppervlakteformule vereist dat $\sigma$ \emph{injectief} is op $K$:
een kaart die een stuk tweemaal overdekt, telt het tweemaal
($\theta$ over $\intcc0{4\pi}$ laten lopen verdubbelt de
\omterm{def:b2:surfaces:area}{oppervlakte} van de bol). (iv) De \omterm{def:b2:surfaces:tangent}{eenheidsnormaal} is door het
oppervlak op het teken na bepaald, maar wordt door de kaart
\emph{gekozen} (de volgorde van $u, v$); uitspraken over oriëntatie
moeten die keuze vastleggen. (v) Ten slotte is $EG - F^2 > 0$ geen
extra hypothese: het is precies de regulariteit, volgens de
identiteit van Lagrange --- verdwijnt zij ergens, dan ligt het
probleem bij de kaart, en geldt daar geen enkele formule voor
\omterm{def:b2:surfaces:area}{oppervlakte} of \omterm{def:b2:surfaces:tangent}{raakvlak}.
\end{remark}

\begin{remark}[Vooruitblik binnen dit volume]
Vooruitwijzende schakels vanaf hier. Het oppervlakte-element
$\sqrt{EG - F^2}\,\dd u\,\dd v$ is een tweedimensionale
jacobideterminant in vermomming, en \cref{ch:b2:multint} maakt de
analogie exact met de stelling over de verandering van
veranderlijken --- de oppervlakte-integralen daar zijn de
\omterm{def:b2:surfaces:area}{oppervlakten} van dit hoofdstuk met een integrand aan boord. De
\omterm{def:b2:surfaces:fff}{eerste fundamentaalvorm} is een veld van positieve \omterm{def:b2:quadratic:def}{kwadratische
vormen}, puntsgewijs behandeld met de gereedschappen van
\cref{ch:b2:quadratic}, waarvan de spectraalstelling ook de
weekendklassering van de \omterm{pb:b2:surfaces:1}{kwadrieken} in dit hoofdstuk aandrijft. En
de normaallijn drijft extremumproblemen op verzamelingen met
nevenvoorwaarden aan (\cref{ex:b2:surfaces:closest}), de
meetkundige kiem van de methode van de multiplicatoren van Lagrange
die bij \cref{thm:b2:diffcalc:extrema} werd geschetst.
\end{remark}

\begin{omfigure}
\begin{tikzpicture}[scale=0.8]
  \draw[omDef, thick] (0,0) ellipse (1.5 and 0.95);
  \draw[omThm] (0,0) ellipse (1.5 and 0.35);
  \node[below, font=\small] at (0,-1.25) {ellipsoïde};
\end{tikzpicture}
\qquad
\begin{tikzpicture}[scale=0.8]
  \draw[omDef, thick] plot[domain=-1.1:1.1, samples=30]
    ({0.6*(exp(\x)+exp(-\x))/2}, {\x});
  \draw[omDef, thick] plot[domain=-1.1:1.1, samples=30]
    ({-0.6*(exp(\x)+exp(-\x))/2}, {\x});
  \draw[omThm] (0,0) ellipse (0.6 and 0.16);
  \draw[omDef] (0,1.1) ellipse (1.0 and 0.22);
  \draw[omDef] (0,-1.1) ellipse (1.0 and 0.22);
  \node[below, font=\small] at (0,-1.55)
    {eenbladige hyperboloïde};
\end{tikzpicture}
\qquad
\begin{tikzpicture}[scale=0.8]
  \draw[omDef, thick] plot[domain=-1.05:1.05, samples=30]
    ({\x}, {1.1*\x*\x - 1.1});
  \draw[omThm] (0,0.055) ellipse (1.0 and 0.22);
  \node[below, font=\small] at (0,-1.45)
    {elliptische paraboloïde};
\end{tikzpicture}
\qquad
\begin{tikzpicture}[scale=0.8]
  \draw[omDef, thick] plot[domain=-1:1, samples=30]
    ({\x}, {0.55*\x*\x});
  \draw[omThm, thick] plot[domain=-1:1, samples=30]
    ({0.45*\x}, {-0.55*\x*\x + 0.18*\x});
  \node[below, font=\small] at (0,-1.15)
    {hyperbolische paraboloïde};
\end{tikzpicture}

{\small Vier van de negen \omterm{pb:b2:surfaces:1}{kwadrieken} die in de weekendopgave worden
geklasseerd, geschetst met hun silhouetten en een niveaukromme
(rood): de begrensde ellipsoïde, de dubbel beschreven eenbladige
hyperboloïde met haar taille, de kom van de elliptische
paraboloïde, en het zadel, waarvan de twee parabolische doorsneden
in tegengestelde richtingen buigen.}
\end{omfigure}

\begin{remark}[Waar dit wordt gebruikt]
Het oppervlakte-element $\norm{\sigma_u\wedge\sigma_v}\,\dd u\,\dd
v$ is de maat van de oppervlakte-integralen van
\cref{ch:b2:multint}, waar het de formule van Green ontmoet; de
\omterm{def:b2:surfaces:fff}{eerste fundamentaalvorm} is het prototype van een veld van
\omterm{def:b2:quadratic:def}{kwadratische vormen}, puntsgewijs bestudeerd met de gereedschappen
van \cref{ch:b2:quadratic}; en de weekendopgave van dit hoofdstuk
klasseert alle \omterm{pb:b2:surfaces:1}{kwadrieken} met de spectraalstelling. Het volume van
bachelorjaar 3 keert tot de oppervlakken terug met
differentiaalvormen en de divergentiestelling, en de intrinsieke
\omterm{thm:b2:curves:frenet2d}{kromming} --- wat $E, F, G$ over het buigen weten --- is de poort
naar de eigenlijke differentiaalmeetkunde.
\end{remark}

\section{Oefeningen}

\begin{exercise}[$\star$]\label{exo:b2:surfaces:1}
Toon aan dat de raakvlakken van de kegel $z = \sqrt{x^2 + y^2}$
(zonder haar top) alle door de top gaan. \emph{(Parametriseer met
$\sigma(\theta, r) = (r\cos\theta, r\sin\theta, r)$, $r > 0$.)}
\end{exercise}

\begin{exercise}[$\star$]\label{exo:b2:surfaces:2}
Bepaal het \omterm{def:b2:surfaces:tangent}{raakvlak} van de ellipsoïde $\frac{x^2}{a^2} +
\frac{y^2}{b^2} + \frac{z^2}{c^2} = 1$ in een punt $(x_0, y_0,
z_0)$ van het oppervlak.
\end{exercise}

\begin{exercise}[$\star$]\label{exo:b2:surfaces:3}
Bereken $E, F, G$ voor de helicoïde $\sigma(u, v) = (v\cos u,\
v\sin u,\ au)$, $a > 0$, en de \omterm{def:b2:surfaces:area}{oppervlakte} van het stuk $0 \leq u
\leq 2\pi$, $0 \leq v \leq 1$, als integraal (evalueer haar met
$\int\sqrt{v^2 + a^2}\,\dd v = \frac12\bigl(v\sqrt{v^2+a^2} +
a^2\ln(v + \sqrt{v^2 + a^2})\bigr) + C$).
\end{exercise}

\begin{exercise}[$\star\star$]\label{exo:b2:surfaces:4}
(Torus) Parametriseer de torus die ontstaat door de cirkel met
middelpunt $(R, 0, 0)$ en straal $r < R$ in het $xz$-vlak om de
$z$-as te draaien:
\[
\sigma(\theta, \psi) = \bigl((R + r\cos\psi)\cos\theta,\
(R + r\cos\psi)\sin\theta,\ r\sin\psi\bigr).
\]
Bereken $E, F, G$, ga de regulariteit na, en toon aan dat de
\omterm{def:b2:surfaces:area}{oppervlakte} $4\pi^2 R r$ is (Pappus: de gemiddelde omtrek $2\pi R$
maal de cirkellengte $2\pi r$).
\end{exercise}

\begin{exercise}[$\star\star$]\label{exo:b2:surfaces:5}
Toon aan dat de \omterm{def:b2:surfaces:area}{oppervlakte} van de grafiek van $f \in
\mathcal{C}^1(K)$ gelijk is aan $\iint_K \sqrt{1 + f_x^2 +
f_y^2}\,\dd x\,\dd y$, en bereken haar voor het stuk paraboloïde $z
= \frac12(x^2 + y^2)$ boven de schijf $x^2 + y^2 \leq 1$
(poolcoördinaten, \cref{ch:b2:multint}).
\end{exercise}

\begin{exercise}[$\star\star$]\label{exo:b2:surfaces:6}
Een getekende kromme $\gamma(t) = \sigma(u(t), v(t))$ op de bol met
straal $R$ (bolkaart) heeft $u = \theta(t)$, $v = \varphi(t)$.
Schrijf haar \omterm{def:b2:curves:length}{lengte} als integraal in $\theta, \varphi$ en bewijs
dat onder de krommen die twee punten van dezelfde meridiaan $\theta
= \theta_0$ verbinden, de meridiaanboog de kortste is. \emph{(Schat
de integrand van onderen af door $R\abs{\varphi'}$.)}
\end{exercise}

\begin{exercise}[$\star\star\star$]\label{exo:b2:surfaces:7}
(Normaallijnen van een bol) Zij $S$ een regulier \omterm{def:b2:metric:connected}{samenhangend}
niveauoppervlak $\{f = c\}$ waarvan alle normaallijnen door een
vast punt $\Omega$ gaan. Toon aan dat $S$ bevat is in een bol met
middelpunt $\Omega$. \emph{(Toon aan dat $\norm{M - \Omega}^2$
afgeleide nul heeft langs elke op $S$ getekende kromme.)}
\end{exercise}

\begin{exercise}[$\star\star\star$]\label{exo:b2:surfaces:8}
(De \omterm{def:b2:surfaces:area}{oppervlakte} is meetkundig) Zij $\Phi \colon U' \to U$ een
$\mathcal{C}^1$-diffeomorfisme tussen \omterm{def:b2:metric:topology}{open} verzamelingen van $\R^2$
en $\tilde\sigma = \sigma \circ \Phi$. Toon aan dat
\[
\norm{\tilde\sigma_{u'} \wedge \tilde\sigma_{v'}}
= \abs{\det J_\Phi}\,
\norm{(\sigma_u \wedge \sigma_v)\circ\Phi} ,
\]
en leid met de formule voor de verandering van veranderlijken uit
\cref{ch:b2:multint} af dat de \omterm{def:b2:surfaces:area}{oppervlakte} uit
\cref{def:b2:surfaces:area} niet van de gekozen reguliere
parametrisering afhangt.
\end{exercise}

\begin{exercise}[$\star$]\label{exo:b2:surfaces:9}
(Hoedendoosstelling van Archimedes) Op de bol met straal $R$ heeft
de \emph{zone} tussen de breedtegraden met $z_1 \leq z \leq z_2$
($-R \leq z_1 < z_2 \leq R$) \omterm{def:b2:surfaces:area}{oppervlakte} $2\pi R\,(z_2 - z_1)$:
bewijs dit met de bolkaart, en besluit dat de \omterm{def:b2:surfaces:area}{oppervlakte} van een
zone alleen van haar hoogte afhangt --- een sinaasappel in
schijven van gelijke dikte snijden geeft even veel schil.
\end{exercise}

\begin{exercise}[$\star\star$]\label{exo:b2:surfaces:10}
Toon aan dat elke normaallijn van een omwentelingsoppervlak
$\sigma(\theta, z) = (r(z)\cos\theta,\ r(z)\sin\theta,\ z)$ (met $r
> 0$ van klasse $\mathcal C^1$) de omwentelingsas snijdt, en
lokaliseer het snijpunt.
\end{exercise}

\begin{exercise}[$\star\star$]\label{exo:b2:surfaces:11}
(De cilinder afrollen) De kaart $\sigma(u, v) = (\cos u,\ \sin u,\
v)$ van de eenheidscilinder heeft $E = G = 1$, $F = 0$: ga dit na
en leg uit waarom elke getekende kromme $t \mapsto \sigma(u(t),
v(t))$ dezelfde \omterm{def:b2:curves:length}{lengte} heeft als de vlakke kromme $t \mapsto
(u(t), v(t))$. Leid af dat de schroeflijn van $(1, 0, 0)$ naar $(1,
0, 2\pi c)$ die één omwenteling maakt, \omterm{def:b2:curves:length}{lengte} $2\pi\sqrt{1 + c^2}$
heeft, en dat geen getekende kromme met dezelfde eindpunten en één
volle omwenteling korter is.
\end{exercise}

\begin{exercise}[$\star\star\star$]\label{exo:b2:surfaces:12}
Zij $S = \{f = c\}$ een \omterm{def:b2:metric:compact}{compact} regulier niveauoppervlak en $M_0
\in S$ een punt op maximale afstand van de oorsprong. Toon aan dat
$\nabla f(M_0)$ collineair is met $\vect{OM_0}$ --- de normaal in
het verste punt is radiaal. Pas dit toe op de ellipsoïde
$\frac{x^2}{a^2} + \frac{y^2}{b^2} + \frac{z^2}{c^2} = 1$ ($a > b >
c > 0$): bepaal alle punten waar de normaal radiaal is, en
identificeer de verste.
\end{exercise}

\section{Probleem: de klassering van de kwadrieken van $\R^3$}

\begin{problem}[{Weekendopgave --- elke kwadriek, gesorteerd door
de spectraalstelling}]\label{pb:b2:surfaces:1}
Een \emph{kwadriek}\index{kwadriek} is de nulverzameling in $\R^3$
van een veelterm van graad twee,
\[
q(X) = X^{\mathsf T}\!AX + 2\,\langle b, X\rangle + c,
\qquad A \in \mathcal S_3(\R),\ A \neq 0,\ b \in \R^3,\
c \in \R .
\]
De oppervlakken van de figuren in dit hoofdstuk --- bollen,
ellipsoïden, zadels, kegels, cilinders --- zijn alle \omterm{pb:b2:surfaces:1}{kwadrieken}.
Deze opgave klasseert ze \emph{alle}: de spectraalstelling
(\cref{thm:b2:quadratic:spectral}) strekt het kwadratische deel,
\omterm{def:b2:affine:subspace}{affiene} translaties (\cref{ch:b2:affine}) slorpen het lineaire deel
op, en wat overblijft is een korte, volledige lijst normaalvormen.

\medskip
\noindent\textbf{Deel I --- De reductiemachine.}
\begin{enumerate}
  \item Zij $X = PY + t$ met $P \in O(3)$ en $t \in \R^3$ (een
        starre coördinatenverandering). Toon aan dat $q(PY + t) =
        Y^{\mathsf T}\!A'Y + 2\langle b', Y\rangle + c'$ met
        \[
        A' = P^{\mathsf T}\!AP, \qquad
        b' = P^{\mathsf T}(At + b), \qquad
        c' = q(t) .
        \]
        Leid af dat het spectrum van $A$ (en dus haar rang en
        signatuur) een invariant onder starre bewegingen van de
        vergelijking is, en leg uit waarom de vergelijking van een
        gegeven \omterm{pb:b2:surfaces:1}{kwadriek} slechts op een scalaire factor ongelijk
        aan nul na bepaald is.
  \item Toon met de spectraalstelling aan dat de vergelijking na
        een rotatie de gedaante $\sum_i \lambda_i y_i^2 +
        2\sum_i\beta_iy_i + c = 0$ aanneemt, met $\lambda_1,
        \lambda_2, \lambda_3$ de \omterm{def:b2:reduction:eigen}{eigenwaarden} van $A$.
  \item Slorp voor elke $i$ met $\lambda_i \neq 0$ de term
        $\beta_iy_i$ op met een translatie ($y_i \mapsto y_i -
        \beta_i/\lambda_i$). Schrijf de herleide vergelijking op
        wanneer $\operatorname{rank} A = r$: $\sum_{i\leq
        r}\lambda_i z_i^2 + 2\sum_{i > r}\beta_iz_i + c'' = 0$.
  \item Een \emph{middelpunt} van de \omterm{pb:b2:surfaces:1}{kwadriek} met vergelijking $q
        = 0$ is een punt $\Omega$ met $q(2\Omega - X) = q(X)$ voor
        alle $X$: de puntspiegeling in $\Omega$ bewaart de
        vergelijking en dus het oppervlak. Toon aan dat $q(2\Omega
        - X) - q(X) = -4\langle A\Omega + b, X\rangle + 4\langle
        A\Omega + b, \Omega\rangle$, en leid af: de middelpunten
        zijn precies de oplossingen van $A\Omega = -b$; zij
        bestaan dan en slechts dan als $b \in
        \operatorname{im}A$, en het middelpunt is uniek dan en
        slechts dan als $A$ inverteerbaar is.
\end{enumerate}

\medskip
\noindent\textbf{Deel II --- Centrale \omterm{pb:b2:surfaces:1}{kwadrieken}
($\operatorname{rank}A = 3$).} Hier is de herleide vergelijking
$\lambda_1z_1^2 + \lambda_2z_2^2 + \lambda_3z_3^2 = \delta$.
\begin{enumerate}[resume]
  \item Neem, zo nodig na vermenigvuldiging met $-1$, aan dat
        minstens twee $\lambda_i > 0$. Som de mogelijkheden op:
        signatuur $(3, 0)$ met $\delta > 0$, $= 0$, $< 0$, en
        signatuur $(2, 1)$ met $\delta > 0$, $= 0$, $< 0$; noem de
        zes resulterende verzamelingen (ellipsoïde, punt, lege
        verzameling, eenbladige hyperboloïde, kegel, tweebladige
        hyperboloïde) en zet elk in haar euclidische normaalvorm
        ($\frac{x^2}{a^2} + \frac{y^2}{b^2} + \frac{z^2}{c^2} =
        1$, enzovoort).
  \item Klasseer $x^2 + y^2 + z^2 + 4xy + 4yz + 4zx = 1$: toon aan
        dat $A = 2J - I$ met $J$ de matrix vol enen, bereken het
        spectrum $\{5, -1, -1\}$, en identificeer een tweebladige
        omwentelingshyperboloïde om de as $\R(1,1,1)$.
  \item (Beschrijvende rechten) Ontbind voor de eenbladige
        hyperboloïde $\frac{x^2}{a^2} + \frac{y^2}{b^2} -
        \frac{z^2}{c^2} = 1$
        \[
        \Bigl(\frac xa - \frac zc\Bigr)
        \Bigl(\frac xa + \frac zc\Bigr)
        = \Bigl(1 - \frac yb\Bigr)\Bigl(1 + \frac yb\Bigr)
        \]
        en breng twee families rechten met één parameter voort die
        op het oppervlak liggen.
  \item Toon aan dat door elk punt van de eenbladige hyperboloïde
        precies één rechte van elke familie gaat: het oppervlak is
        \emph{dubbel beschreven}.
  \item De \emph{asymptotische kegel} van de eenbladige
        hyperboloïde is $C : \frac{x^2}{a^2} + \frac{y^2}{b^2} -
        \frac{z^2}{c^2} = 0$. Toon met $\rho^2 = \frac{x^2}{a^2} +
        \frac{y^2}{b^2}$ aan dat elk punt van de hyperboloïde op
        afstand hoogstens $c\bigl(\rho - \sqrt{\rho^2 - 1}\bigr) =
        \frac{c}{\rho + \sqrt{\rho^2 - 1}}$ van $C$ ligt, zodat
        het oppervlak zich in het oneindige tegen haar kegel
        aanvlijt. Wat zijn de doorsneden van de hyperboloïde met
        de vlakken $x = \pm a$?
\end{enumerate}

\medskip
\noindent\textbf{Deel III --- Rang $2$ en rang $1$:
paraboloïden, cilinders, vlakken.}
\begin{enumerate}[resume]
  \item Neem aan dat $\operatorname{rank}A = 2$, zeg $\lambda_1,
        \lambda_2 \neq 0 = \lambda_3$. Splits, vertrekkend van
        vraag 3, in twee gevallen naargelang $\beta_3 \neq 0$
        (geen middelpunt, volgens vraag 4) of $\beta_3 = 0$ (een
        rechte middelpunten), en herleid tot
        \[
        \lambda_1z_1^2 + \lambda_2z_2^2 + 2\beta_3z_3 = 0
        \qquad\text{of}\qquad
        \lambda_1z_1^2 + \lambda_2z_2^2 + c'' = 0 :
        \]
        elliptische of hyperbolische paraboloïden in het eerste
        geval, cilinders boven centrale kegelsneden (of paren
        snijdende vlakken, een rechte, de lege verzameling) in het
        tweede.
  \item Toon aan dat het zadel $z = xy$ een hyperbolische
        paraboloïde is: draai over $\pi/4$ in het $xy$-vlak om $z
        = \tfrac12(u^2 - v^2)$ te bereiken, op de schaal na het
        oppervlak van \cref{fig:b2:surfaces:saddle}.
  \item Toon aan dat het zadel $z = xy$ de twee families rechten
        $\{x = x_0,\ z = x_0y\}$ en $\{y = y_0,\ z = xy_0\}$
        draagt, met door elk punt precies één rechte van elke
        familie: de tweede dubbel beschreven \omterm{pb:b2:surfaces:1}{kwadriek}.
  \item Klasseer $x^2 + y^2 - 2x + 4y + 3 = 0$ in $\R^3$ (maak de
        kwadraten af; identificeer een rechte cirkelvormige
        cilinder, en geef haar as en straal).
  \item Zij nu $\operatorname{rank}A = 1$, zeg $\lambda_1 \neq 0 =
        \lambda_2 = \lambda_3$. Herleid, door binnen het kernvlak
        te draaien en te transleren, tot
        \[
        \lambda_1z_1^2 + 2\beta z_2 = 0 \quad (\beta \neq 0)
        \qquad\text{of}\qquad
        \lambda_1z_1^2 + c'' = 0 :
        \]
        een parabolische cilinder, of een paar evenwijdige
        vlakken, een dubbel vlak, of de lege verzameling.
        Klasseer $(x + y)^2 = z$ \omterm{def:b2:metric:complete}{volledig} (normaalvorm, as van
        translatie-invariantie).
\end{enumerate}

\medskip
\noindent\textbf{Deel IV --- De klasseringsstelling.}
\begin{enumerate}[resume]
  \item Zet de Delen I--III in elkaar tot een stelling: \emph{elke
        \omterm{pb:b2:surfaces:1}{kwadriek} van $\R^3$ wordt door een starre beweging op
        precies één normaalvorm afgebeeld}. Som de zeventien
        \omterm{def:b2:affine:subspace}{affiene} types op (tel de lege varianten en de ontaarde
        verzamelingen mee), en zonder de negen \omterm{pb:b2:surfaces:1}{kwadrieken} die
        \emph{oppervlakken} zijn af: ellipsoïde, een- en
        tweebladige hyperboloïde, kegel, elliptische en
        hyperbolische paraboloïde, en elliptische, hyperbolische
        en parabolische cilinder.
  \item Schrijf het \emph{klasseringsalgoritme} op: welke
        grootheden bereken je, gegeven $(A, b, c)$, in welke
        volgorde, en welke vertakking beslist over welk type?
        Verantwoord dat elke stap uitvoerbaar is (\omterm{def:b2:reduction:eigen}{eigenwaarden} van
        een \omterm{def:b2:quadratic:adjoint}{symmetrische} $3\times3$-matrix, rang, oplosbaarheid
        van $A\Omega = -b$).
  \item Laat het algoritme lopen op $x^2 + y^2 + z^2 - 2xy - 2yz -
        2zx = 1$: toon aan dat $A = 2I - J$ spectrum $\{2, 2,
        -1\}$ heeft en besluit: een eenbladige
        omwentelingshyperboloïde om $\R(1,1,1)$.
  \item Laat het lopen op $x^2 + y^2 - z^2 - 2x + 4y + 2z + 4 =
        0$: bepaal het middelpunt en identificeer de \omterm{pb:b2:surfaces:1}{kwadriek}.
  \item Laat het lopen op $x^2 + 4xy + y^2 = 2z$: diagonaliseer
        het $xy$-blok ($u = \frac{x+y}{\sqrt2}$, $v =
        \frac{x-y}{\sqrt2}$) en identificeer de \omterm{pb:b2:surfaces:1}{kwadriek}.
  \item Euclidisch tegenover affien. Toon aan dat twee centrale
        \omterm{pb:b2:surfaces:1}{kwadrieken} in normaalvorm \emph{star} gelijkwaardig zijn
        dan en slechts dan als zij dezelfde lijsten coëfficiënten
        hebben (op een permutatie en een gemeenschappelijke
        positieve schaalfactor op de vergelijking na), terwijl er
        \emph{affien} alleen de gegevens van de signatuur
        overleven: elke ellipsoïde is een affien beeld van de
        ronde bol. Welke stelling waarborgt dat de signatuur
        onderweg niet kan veranderen
        (\cref{thm:b2:quadratic:sylvester})?
\end{enumerate}

\medskip
\noindent\textbf{Deel V --- Dividenden.}
\begin{enumerate}[resume]
  \item Toon aan dat elke doorsnede van een \omterm{pb:b2:surfaces:1}{kwadriek} met een
        affien vlak een (mogelijk ontaarde) kegelsnede van dat
        vlak is. Identificeer de doorsneden van het zadel $z = xy$
        met de vlakken $z = c$ ($c \neq 0$ en $c = 0$).
  \item Welke \omterm{pb:b2:surfaces:1}{kwadrieken} bevatten rechten? Toon aan dat de
        ellipsoïde, de tweebladige hyperboloïde en de elliptische
        paraboloïde er geen bevatten \emph{(beperk $q$ tot een
        rechte en gebruik de ongelijkheid van Cauchy--Schwarz voor
        het tweebladige geval)}; dat kegel en cilinders door één
        familie worden beschreven; en besluit dat de dubbel
        beschreven \omterm{pb:b2:surfaces:1}{kwadrieken} precies de eenbladige hyperboloïde
        en de hyperbolische paraboloïde zijn.
  \item Wil men alleen het \emph{\omterm{def:b2:affine:subspace}{affiene}} type, dan is de reductie
        van Gauss (\cref{thm:b2:quadratic:gauss}) goedkoper dan
        diagonaliseren. Doe vraag 17 opnieuw met het algoritme van
        Gauss en ga de signatuur $(2, 1)$ na; welke euclidische
        informatie verliest Gauss?
  \item Alle \omterm{def:b2:reduction:eigen}{eigenwaarden} van een \omterm{def:b2:quadratic:adjoint}{symmetrische} matrix zijn reëel;
        toon aan dat de \emph{tekens} van de \omterm{def:b2:reduction:eigen}{eigenwaarden} van $A$
        bijgevolg met de tekenregel van Descartes uit de
        \omterm{def:b2:reduction:charpoly}{karakteristieke veelterm} kunnen worden afgelezen, en ga
        het na op vraag 6: $\chi_A(\lambda) = \lambda^3 -
        3\lambda^2 - 9\lambda - 5$ heeft precies één
        tekenwisseling, en dus signatuur $(1, 2)$.
  \item Synthese. Vat het algoritme in enkele regels samen;
        formuleer de exacte rol van (i) de spectraalstelling, (ii)
        de middelpuntsvergelijking $A\Omega = -b$, (iii) de
        traagheidsstelling van Sylvester, (iv) de reductie van
        Gauss. Wat levert dezelfde machine in $\R^2$ op, en wat
        verandert er in $\R^n$?
\end{enumerate}
\end{problem}
